\section{萨克斯--乌伦贝克重标度与冒泡}
\stfhIndex{I240}{sa^ke^si^wu^lun^bei^ke^chong^biao^du}{萨克斯--乌伦贝克重标度（Sacks–Uhlenbeck rescaling）}
\label{sec:sacks-uhlenbeck-rescaling-bubbling}

现在研究固定同调类中一列 $J$-全纯映射 $(j_\alpha,u_\alpha)$ 的收敛性质，特别假设其调和能量有一致上界。需要强调，定理 8.3.5 对 $k\geq2$ 建立的主要 $W^{k,p}$ 强制估计依赖于 $p>2$ 时的 $W^{1,p}$-范数。因此，关键问题是在有限能量假设，即映射导数的 $L^2$ 上界下，建立 $W^{1,p}$ 估计。在更高正则性估计已经具备时，这等价于：取子列并对映射作适当重参数化后，为给定序列建立一致导数界。

从 $C^1$ 估计与阿斯科利--阿尔泽拉定理所给出的这一观点看，格罗莫夫紧性定理（特别是定义于亏格 $0$ 的黎曼曲面上的映射）可以改述为：去掉有限个“气泡”后，存在一致导数界。每个气泡都是一个 $J$-全纯球面（或圆盘），并消耗不小于某个一致正下界的非零能量。构造这种气泡时，先在平面
$\mathbb C=\mathbb{CP}^1\setminus\{\infty\}$ 上采用著名的萨克斯--乌伦贝克重标度论证，再应用上一章给出的可去奇点定理。

为了突出萨克斯--乌伦贝克重标度过程，本节主要关注闭黎曼曲面 $(\Sigma,j)$，并固定 $\Sigma$ 上的共形结构 $j$。下一节将在稳定映射的背景下完整讨论格罗莫夫紧性。在这种情况下，把 $(j,J)$-全纯映射简称为 $J$-全纯映射。我们在证明 $\epsilon$-正则性定理，即定理 8.4.1 时，已经使用过萨克斯--乌伦贝克重标度论证最基本的版本。

假设 $M$ 紧且无边界；或者当 $M$ 非紧时，假设 $(M,\omega,J)$ 在定义 8.7.7 的意义下是驯服的。还固定定义域复结构 $j$，从而固定 $\Sigma$ 上的度量 $h$。

设 $u_\alpha:(\Sigma,j)\to(M,J)$ 是一列满足下述条件的 $J$-全纯映射：
\begin{enumerate}[label=(\arabic*)]
 \item $E(u_\alpha)<C$；
 \item 存在紧子集 $K\subset M$，使得对所有 $\alpha$ 都有
 $\operatorname{Image}u_\alpha\subset K$。
\end{enumerate}
则下列两种情形之一成立：
\begin{enumerate}[label=(\arabic*)]
 \item 存在 $u_{\alpha_j}$ 的一个子列，使得可以选取 $\Sigma$ 的一列双全纯映射 $\phi_j$，满足
 \[
  \max_{z\in\Sigma}|d(u_{\alpha_j}\circ\phi_j)(z)|\leq C_1
 \]
 对某个 $C_1>0$ 成立；
 \item 对任意一列双全纯映射 $\phi_\alpha$，复合映射
 $u_\alpha\circ\phi_\alpha$ 都没有子列
 $u_{\alpha_j}\circ\phi_{\alpha_j}$，使得
 $\|d(u_{\alpha_j}\circ\phi_{\alpha_j})\|_{C^0}$ 一致有界。
\end{enumerate}

先考虑情形 (1)。把 $u_\alpha\circ\phi_\alpha$ 仍简记为 $u_\alpha$，并假设存在 $C_1>0$，使所有 $\alpha$ 都满足
$\|du_\alpha\|_{C^0}\leq C_1$。由于定义域 $\Sigma$ 紧，这蕴含
\[
 \|du_\alpha\|_p\leq C_2(p,C_1)
\]
对任意 $p>0$ 成立。固定定理 8.3.5 中所取的某个 $p>2$。把主要估计定理 8.3.5 应用于 $\Sigma$ 的覆盖 $\{D_\beta\}$ 的每个元素，得到
$\|du_\alpha\|_{1,p}$ 的界。再由索伯列夫嵌入
\stfhIndex{I252}{suo^bo^lie^fu^qian^ru}{索伯列夫嵌入（Sobolev embedding）}
$W^{1,p}\hookrightarrow C^\epsilon$（$0\leq\epsilon<1-2/p$），可以抽取 $u_\alpha$ 的一个子列，它在 $C^\epsilon$-范数中收敛于映射 $u_0\in C^\epsilon$。这个极限映射 $u_0$ 是位于 $C^\epsilon$ 中的 $\overline\partial_Ju=0$ 的弱解。主要正则性定理 8.5.5 随即蕴含 $u_0$ 实际上光滑；通过自举论证，收敛成为 $C^\infty$-收敛。特别地，极限映射 $u_0$ 与
$[u_\alpha]=\alpha\in H_2(M)$ 具有相同同调类。至此结束对情形 (1) 的讨论。

下面的定理描述一般 $J$-全纯映射序列 $u_\alpha$ 在情形 (2) 中会发生什么，即
\begin{equation}
 E(u_\alpha)<C,
 \qquad \max_{z\in S^2}|du_\alpha(z)|\longrightarrow\infty.
 \label{eq:bubbling-energy-derivative-assumption}
\end{equation}
回忆命题 8.4.8，下面的常数严格为正：
\[
 A=A(\omega;J)
 =\inf_u\left\{\int u^*\omega\mathrel{\Big|}
 \overline\partial_Ju=0, u:S^2\to M, u\text{ 非常值。}\right\}>0.
\]

\begin{theorem}\label{thm:bubbling-measure-convergence}
设 $u_\alpha:\Sigma\to M$ 如\eqref{eq:bubbling-energy-derivative-assumption}。则存在一个仍记为 $\{u_\alpha\}$ 的子列、某个 $L\in\mathbb Z_+$ 对应的有限点集
$\{x_1,\ldots,x_L\}\subset\Sigma$，以及一个 $J$-全纯映射
$u_0:\Sigma\to M$，使得：
\begin{enumerate}[label=(\arabic*)]
 \item $u_\alpha\to u_0$ 在
 $\Sigma\setminus\{x_1,\ldots,x_L\}$ 上按紧 $C^1$ 拓扑收敛；
 \item 把能量密度 $e(u_\alpha)=|du_\alpha|^2$ 看作测度
 $e(u_\alpha)d\mu$（$\mu$ 是与 $(\Sigma,h)$ 相伴的测度），它收敛于 $e(u_0)d\mu$ 加上一列点测度之和：
 \[
  e(u_\alpha)\longrightarrow e(u_0)+\sum_{i=1}^Lm_i\delta_{x_i},
 \]
 其中每个 $m_i\geq A(\omega;J)>0$。
\end{enumerate}
\end{theorem}

\begin{proof}
固定 $r_0>0$，并考虑序列 $r_m=2^{-m}r_0$，其中
$m\in\mathbb Z_+$；选取 $r_0$ 小于 $(\Sigma,h)$ 的单射半径。

对每个 $m\in\mathbb Z_+$，选取 $\Sigma$ 的有限覆盖
\[
 \mathcal U_m=\{D_{r_m}(y_i)\mid y_i\in\Sigma\},
\]
使得 $D_{r_m/2}(y_i)$ 也覆盖 $\Sigma$，而每一点 $x\in\Sigma$ 至多被 $\mathcal U_m$ 中的球覆盖 $l$ 次；这里 $l$ 可以与 $m$ 无关，例如因为 $\Sigma$ 是二维曲面，取 $l=3$ 即已足够。

由假设 $\int_\Sigma|du_\alpha|^2<C$，有
\[
 \sum_{i=1}^l\int_{D_{r_m}(y_i)}|du_\alpha|^2\leq l\cdot C.
\]
因此，满足
$\int_{D_{r_m}(y_i)}|du_\alpha|^2\geq\epsilon_0$ 的圆盘至多有
$[lC/\epsilon_0]+1$ 个。（这里把 $\epsilon_0$ 取为主要密度估计中的常数；见定理 8.4.1 或命题 8.4.6。）称这种圆盘为“坏”圆盘。

由于 $lC/\epsilon_0$ 与 $\alpha$ 无关，而 $\Sigma$ 紧，当
$\alpha\to\infty$ 时，这些坏圆盘的中心构成 $\Sigma$ 中的点列。用
\[
 \{x_{1,m}^\alpha,\ldots,x_{L,m}^\alpha\}
\]
表示这些坏圆盘的中心，其中整数 $L\leq[lC/\epsilon_0]+1$。取 $\alpha$ 的子列后，可以假设当 $\alpha\to\infty$ 时它们收敛于
\[
 \{x_{1,m},\ldots,x_{L,m}\}.
\]
再取 $m$ 的一个子列，还可假设对所有 $i=1,\ldots,L$，当
$m\to\infty$ 时 $x_{i,m}\to x_i$；下文都作此假设。

除这些坏圆盘外，在每个好圆盘 $D_{r_m}(y_i)$ 上都有
\[
 \int_{D_{r_m}(y_i)}|du_\alpha|^2\leq\epsilon_0.
\]
在这些好圆盘上，可以应用密度估计，得到 $D_{r_m/2}(y_i)$ 上
$du_\alpha(x)$ 的界。因此，问题归结为先前考虑的第一种情形；取 $u_\alpha$ 的一个子列后，$u_\alpha$ 在 $D_{r_m/2}(y_i)$ 上收敛。

注意，对任意给定的紧子集
$K\subset\Sigma\setminus\{x_1,\ldots,x_L\}$，当 $m,\alpha$ 充分大时都有
\[
 K\subset\Sigma-\bigcup_{k=1}^LD_{r_m/2}(x_{k,m}^\alpha).
\]
对每个固定的 $m$ 先令 $\alpha\to\infty$，再在 $m\to\infty$ 时令 $K$ 穷竭
$\Sigma\setminus\{x_1,\ldots,x_L\}$。取 $\alpha=\alpha(m)$ 在
$m\to\infty$ 时的一个对角子列，得到极限
$u_\alpha\to u_0$，它在
$\Sigma\setminus\{x_1,x_2,\ldots,x_L\}$ 上按紧 $C^1$ 拓扑（从而按紧 $C^\infty$ 拓扑）收敛。由于
\[
 \int_{\Sigma\setminus\{x_1,\ldots,x_L\}}|du_0|^2
 \leq\liminf_{\alpha\to\infty}\int|du_\alpha|^2\leq E_0
\]
且 $u_0$ 在 $\Sigma\setminus\{x_1,\ldots,x_L\}$ 上满足
$\overline\partial_Ju_0=0$，可去奇点定理蕴含 $u_0$ 越过
$\{x_1,\ldots,x_L\}$ 光滑延拓到整个定义域 $\Sigma$。这证明了陈述 (1)。

现在证明定理的陈述 (2)。已经找到了 $C^1$-收敛失效的点
$x_1,x_2,\ldots,x_L$。对任意 $\epsilon>0$，数
\[
 b_\alpha^k=\sup\{|du_\alpha(x)|\mid x\in D_\epsilon(x_k)\}
\]
在 $\alpha\to\infty$ 时必无界。取子列后，可以假设对每个
$k=1,\ldots,L$，都有 $b_\alpha^k\to\infty$。固定 $\epsilon>0$，使其小于单射半径的一半，并且小于
\[
 \frac14\min_{i,j}\{d(x_i,x_j)\mid1\leq i\ne j\leq L\};
\]
于是 $D_{2\epsilon}(x_i)$ 是嵌入且两两不交的。

令
\[
 m_k:=\lim_{\epsilon\to0}\limsup_{\alpha\to\infty}
 \int_{D_\epsilon(x_k)}
 \bigl||du_\alpha|^2-|du_0|^2\bigr|\,dA.
\]
需要对任意 $1\leq k\leq L$ 证明 $m_k\geq A(\omega,J)$。在下面的讨论中固定 $k$，从而把 $x_k$ 简记为 $x$，把 $b_\alpha^k$ 简记为 $b_\alpha$。以 $\overline x_\alpha$ 表示
$D_\epsilon(x)$ 中满足
\[
 |du_\alpha(\overline x_\alpha)|
 =\max_{y\in D_\epsilon(x)}|du_\alpha(y)|=b_\alpha
\]
的点。

\begin{lemma}\label{lem:bubbling-maximum-double-disc}
给定 $\epsilon>0$。对所有充分大的 $\alpha$，有
\begin{equation}
 b_\alpha=|du_\alpha(\overline x_\alpha)|
 =\max_{y\in D_{2\epsilon}(x)}|du_\alpha(y)|.
 \label{eq:bubbling-maximum-double-disc}
\end{equation}
\end{lemma}

\begin{proof}
由 $u_\alpha$ 在
$\Sigma-\bigcup_{k=1}^LD_\epsilon(x_{k,m}^\alpha)$ 上按紧 $C^1$ 拓扑收敛，有
\[
 \max_{y\in\Sigma-\bigcup_{k=1}^LD_\epsilon(x_k)}|du_\alpha(y)|
 \leq C(\epsilon)<\infty.
\]
特别地，
\[
 \max_{y\in D_{2\epsilon}(x_k)\setminus D_\epsilon(x_k)}
 |du_\alpha(y)|\leq C(\epsilon).
\]
由于 $b_\alpha\to\infty$，最终有 $b_\alpha>C(\epsilon)$，证明完毕。
\end{proof}

于是，由上述设置，$\alpha\to\infty$ 时
$\overline x_\alpha\to x$，特别是
$D_\epsilon(\overline x_\alpha)\subset D_{2\epsilon}(x)$。在
$D_{2\epsilon}(x)\subset\Sigma$ 上固定坐标系，并假设
$D_{2\epsilon}(x)\subset\mathbb C$。现在重归一化 $u_\alpha$，在
$D_{\epsilon b_\alpha}(0)$ 上定义映射
\[
 \widetilde u_\alpha(z)
 =u_\alpha\left(\overline x_\alpha+\frac z{b_\alpha}\right).
\]
于是 $|d\widetilde u_\alpha(0)|=1$。注意，对所有
$z\in D_{\epsilon b_\alpha}(0)$，都有
\[
 \overline x_\alpha+\frac z{b_\alpha}\in D_{2\epsilon}(x).
\]
这与\eqref{eq:bubbling-maximum-double-disc}蕴含
\[
 |d\widetilde u_\alpha|(z)
 =\frac1{b_\alpha}
 |du_\alpha(\overline x_\alpha+z/b_\alpha)|
 \leq\frac1{b_\alpha}
 \max_{y\in D_{2\epsilon}(x)}|du_\alpha(y)|
 =\frac1{b_\alpha}b_\alpha=1
\]
对所有 $z\in D_{\epsilon b_\alpha}(0)$ 成立。此外，当
$\alpha\to\infty$ 时，$D_{\epsilon b_\alpha}(0)$ 穷竭整个
$\mathbb C$，并满足
\[
 \int_{D_{\epsilon b_\alpha}(0)}|d\widetilde u_\alpha|^2\leq E_0.
\]

现在固定 $R>0$，并在 $D_R(0)$ 上考虑 $\widetilde u_\alpha$。当
$\alpha\to\infty$ 时，最终有
$D_R(0)\subset D_{\epsilon b_\alpha}(0)$。由于在 $D_R(0)$ 上
$|d\widetilde u_\alpha|\leq1$，存在 $\widetilde u_\alpha$ 的一个子列，记为 $\widetilde u_{\alpha,R}$，使得当 $\alpha\to\infty$ 时，
$\widetilde u_{\alpha,R}$ 在 $D_{R-1}(0)$ 上按 $C^1$ 拓扑收敛。以
$u_R:D_{R-1}(0)\to M$ 表示相应极限。由
$\widetilde u_{\alpha,R}\to u_R$ 的 $C^1$-收敛与
$|d\widetilde u_\alpha(0)|=1$，有 $|du_R(0)|=1$。

可以取 $R=N\in\mathbb N$。令 $N\to\infty$，并选取对角子列
$\alpha=\alpha(N)$，使得
\[
 \{\widetilde u_{\alpha,N+1}\}\subset
 \{\widetilde u_{\alpha,N}\}.
\]
因此，对角子列有定义在 $\mathbb C$ 上的极限 $u_\infty$。由紧
$C^1$-收敛，$u_\infty$ 满足
$|du_\infty(0)|=1$、$|du_\infty|\leq1$，并且
$\widetilde u_{\alpha,N}$ 逐点收敛于 $u_\infty$，从而
$|d\widetilde u_{\alpha,N}|^2\to|du_\infty|^2$ 几乎处处成立。由法图引理，
\[
 \int_{\mathbb C}|du_\infty|^2
 \leq\liminf_N\int_{\mathbb C}|d\widetilde u_{\alpha,N}|^2
 \leq E_0.
\]

另一方面，固定共形等同
$\mathbb C\cong S^2\setminus\{p\}$，并把上述所有映射以及
$u_\infty$ 看成定义于 $S^2\setminus\{p\}$ 的子集上的映射。由可去奇点定理，$u_\infty$ 延拓到 $S^2$。由于
$|du_\infty(0)|=1$，$u_\infty$ 不是常值映射，所以依 $A(\omega,J)$ 的定义，
\[
 E(u_\infty)\geq A(\omega,J).
\]
恢复指标 $k=1,\ldots,L$，并注意当 $\epsilon\to0$ 时
$\int_{D_\epsilon(x_i)}|du_0|^2dA\to0$，由 $u_\infty$ 的构造得到
\[
\begin{split}
 E(u_\infty)
 &=\int_{\mathbb C}|du_\infty|^2
 \leq\liminf_N\int_{D_N}|d\widetilde u_{\alpha,N}|^2\\
 &\leq\lim_{\epsilon\to0}\limsup_{\alpha\to\infty}
 \int_{D_\epsilon(x_k)}
 \bigl||du_\alpha|^2-|du_0|^2\bigr|\,dA
 =m_k.
\end{split}
\]
这就完成了关于测度 $e(u_\alpha)$ 收敛的陈述 (2) 之证明。
\end{proof}

\begin{remark}\label{rem:bubbling-boundary-principal-components}
\begin{enumerate}[label=(\arabic*)]
 \item 上述定理同样适用于有边界 $\partial\Sigma$ 的情形。唯一差别是要加入圆盘气泡映射
 $w_\infty:(D^2,\partial D^2)\to(M,L)$。（见上一章 $\epsilon$-正则性定理的证明。）
 \stfhIndex{I230}{zhu^fen^zhi}{主分支（principal component）}
 \item 主分支 $u_0$ 可以是常值映射。（注意，依定义，气泡不能是常值映射。）用圆盘映射的情形说明这个现象。考虑图 9.1 所示的共形映射。

 这里把 $\mathbb R\times[0,1]$ 与 $D^2\setminus\{\pm1\}$ 等同，并以 $u_r$ 表示 $-1<r<1$ 时的全纯映射，它把
 $(0,0)\mapsto(r,0)$、$(0,1)\mapsto(1,0)$。当 $r\to1-0$ 时，
 $u_r$ 的导数在 $(0,1)$ 处爆破；这产生一个附着于点
 $(0,-1)\in\mathbb C$ 其他点处的圆盘气泡 $u_\infty$，而主部
 $u_0$ 成为常值映射。
 \item 在定理 9.2.1 的证明中，构造了两个 $J$-全纯映射
 $u_0$ 与 $u_\infty$。虽然 $u_\infty$ 贡献 $\delta$ 质量
 $m_i\delta_{x_i}$，但 $u_\infty$ 的像不一定附着在点
 $u_0(x_i)$。如果在主分支 $u_0$ 与气泡分支 $u_\infty$ 之间的环形区域上还出现另一个气泡，就可能发生这种情况。
\end{enumerate}
\end{remark}

\begin{figure}[htbp]
 \centering
 \begin{tikzpicture}[x=1cm,y=1cm,>=Stealth,line cap=round,line join=round]
  \begin{scope}[shift={(-3.2,0)}]
    \fill[gray!35] (-1.45,-.42) rectangle (1.45,.42);
    \draw[densely dotted] (-1.45,-.42) rectangle (1.45,.42);
    \draw[->] (-1.55,0)--(1.58,0);
    \draw[->] (0,-1.15)--(0,1.2);
    \node[below right] at (0,0) {$0$};
    \node[above right] at (0,.42) {$i$};
  \end{scope}
  \draw[->,decorate,decoration={snake,amplitude=1.3mm,segment length=5mm}]
    (-1.25,0)--(.25,0);
  \begin{scope}[shift={(2.15,0)}]
    \fill[gray!35] (0,0) circle (.8);
    \draw (0,0) circle (.8);
    \draw[->] (-1.25,0)--(1.35,0);
    \draw[->] (0,-1.15)--(0,1.2);
    \draw[fill=white] (-.8,-.055) rectangle (.18,.055);
  \end{scope}
\end{tikzpicture}

 \caption{共形映射。}
 \label{fig:conformal-map-strip-disc}
\end{figure}

从前几节对能量界下紧性失效的分析已经看到，对 $J$-全纯曲线集合作合适紧化时应当包含奇异曲线。

格罗莫夫引入了包含“尖点曲线”的伪全纯曲线紧化 (Gr85)。这个结果通常称为格罗莫夫紧性定理。随后，皮埃尔·潘苏（Pierre Pansu）(Pa94)、帕克与沃尔夫森 (PW93) 为闭情形，叶如钢 (Ye94) 为一般情形，进一步奠定了紧性定理的证明基础。孔采维奇引入了\emph{稳定映射}概念，从而给出一个如今广泛使用的伪全纯映射模空间紧化。特别地，这种紧化导出一个\emph{紧且豪斯多夫}的拓扑。它还容许良好的无穷小形变理论；在代数几何中表现为\emph{完美障碍理论} (LiT98)，在光滑背景中则表现为仓西结构 (FOn99)（或其衍生形式）。
